\chapter{Sistemas Proteicos}
\label{cap:07}

Se o genoma é o manual de instruções da célula, as proteínas são as máquinas que o
executam. Elas catalisam reações que de outro modo levariam séculos, transmitem sinais
do ambiente ao núcleo, constroem o citoesqueleto, transportam oxigênio, reconhecem
patógenos e regulam a própria expressão dos genes. Compreender os sistemas proteicos ---
sua estrutura tridimensional, sua cinética, as cascatas de sinalização que coordenam,
as modificações pós-traducionais que ampliam sua diversidade funcional e as tecnologias
que permitem medi-los em escala genômica --- é indispensável para qualquer abordagem
quantitativa da biologia de sistemas.

Este capítulo organiza-se em cinco temas: estrutura de proteínas, cinética enzimática
avançada, cascatas de sinalização celular, modificações pós-traducionais e proteômica.
Em cada tema, conectamos a descrição molecular com modelos matemáticos e ferramentas
computacionais.

\section{Estrutura de Proteínas}

\subsection{Níveis de Organização Estrutural}

A estrutura de uma proteína é descrita em quatro níveis hierárquicos. A \emph{estrutura
primária} é a sequência linear de aminoácidos ligados por ligações peptídicas --- a
informação diretamente codificada pelo mRNA. A \emph{estrutura secundária} consiste em
arranjos locais regulares estabilizados por ligações de hidrogênio entre átomos da cadeia
principal: a \emph{$\alpha$-hélice} (3.6 resíduos por volta, $\sim$5.4 Å de passo) e a
\emph{$\beta$-folha} (paralela ou antiparalela). A \emph{estrutura terciária} é o arranjo
tridimensional completo de uma cadeia polipeptídica única, determinado pela interação entre
cadeias laterais. A \emph{estrutura quaternária} descreve o arranjo de múltiplas subunidades
polipeptídicas num complexo.

O dobramento proteico é termodinamicamente governado pela minimização da energia livre de
Gibbs $\Delta G = \Delta H - T\Delta S$. O \emph{efeito hidrofóbico} é a força motriz
dominante: resíduos apolares, expostos à água no estado desnaturado, são energeticamente
desfavoráveis e buscam o interior da proteína ao dobrar, aumentando a entropia do solvente
(os moléculas de água que os circundavam são liberadas). As ligações de hidrogênio, as
interações eletrostáticas e as forças de Van der Waals contribuem para $\Delta H$; a perda
de liberdade conformacional do polipeptídeo contribui negativamente para $\Delta S$. A
margem de estabilidade termodinâmica da maioria das proteínas é de apenas 5--20 kcal/mol
--- o equivalente a poucas ligações de hidrogênio.

O \emph{diagrama de Ramachandran} organiza as conformações permitidas da cadeia principal
em função dos ângulos de torção $\phi$ (N--C$_\alpha$) e $\psi$ (C$_\alpha$--C). Regiões
favoráveis correspondem a $\alpha$-hélices ($\phi \approx -60^\circ$, $\psi \approx -45^\circ$)
e $\beta$-folhas ($\phi \approx -120^\circ$, $\psi \approx +120^\circ$). Regiões proibidas
ocorrem onde cadeias laterais colidem esfericamente com átomos da cadeia principal.
Glicina, sem cadeia lateral, tem liberdade conformacional em praticamente toda a área do
diagrama; prolina, com nitrogênio imino em anel, restringe $\phi$ a $\approx -60^\circ$.

\subsection{Domínios, Motivos e Proteínas Desordenadas}

Os \emph{domínios proteicos} são unidades estruturais e funcionais independentes ---
geralmente 50--200 aminoácidos --- capazes de dobrar-se autonomamente. A arquitetura
multidomínio permite que proteínas combinem funções modularmente: um domínio quinase
com um domínio de regulação, um domínio de ligação ao DNA com um domínio de ativação
transcricional. A evolução molecular opera amplamente por recombinação e duplicação de
domínios, gerando a diversidade do proteoma sem criar cada função do zero.

Os \emph{motivos estruturais} são padrões recorrentes em nível de estrutura secundária:
hélice-volta-hélice (reconhecimento de DNA), dedos de zinco (fator de transcrição),
barril TIM (enzimas metabólicas), leucine zipper (dimerização). Esses motivos são unidades
evolutivas conservadas que aparecem em proteínas funcionalmente diversas.

As \emph{proteínas intrinsecamente desordenadas} (IDPs) desafiam o paradigma clássico
estrutura--função: 30--40\% do proteoma eucariótico contém regiões sem estrutura 3D
definida em condições fisiológicas. Ricas em resíduos polares e carregados e pobres em
hidrofóbicos, as IDPs exibem flexibilidade conformacional que lhes confere vantagens
funcionais: podem interagir com múltiplos parceiros diferentes (promiscuidade de ligação),
são altamente reguláveis por modificações pós-traducionais e permitem montagem de complexos
de sinalização transitórios. O supressor tumoral p53, a proteína $\alpha$-sinucleína
(central na doença de Parkinson) e as caudas N-terminais das histonas são exemplos
biologicamente centrais.

\subsection{O Protein Data Bank}

O Protein Data Bank (PDB, rcsb.org) é o repositório global de estruturas tridimensionais
de macromoléculas biológicas, com mais de 200.000 estruturas determinadas por
cristalografia de raios-X (85\%), ressonância magnética nuclear (NMR) e microscopia
crioeletrônica (cryo-EM). Cada estrutura é identificada por um código de 4 caracteres e
armazenada em formato PDB ou mmCIF, contendo coordenadas atômicas (ATOM/HETATM), metadados
experimentais e anotação de estrutura secundária (registros HELIX, SHEET).

\begin{lstlisting}[language=Python, caption=Analise de estruturas PDB com BioPython]
from Bio.PDB import PDBParser, DSSP
import numpy as np

parser = PDBParser(QUIET=True)
structure = parser.get_structure("1UBQ", "1ubq.pdb")
model = structure[0]

# Calcular centro de massa da cadeia A
chain = model["A"]
coords = np.array([atom.coord for res in chain
                   for atom in res if res.id[0] == ' '])
print(f"Centro de massa: {coords.mean(axis=0)}")
print(f"Raio de giro: {np.sqrt(((coords - coords.mean(0))**2).sum(1).mean()):.2f} A")

# Estrutura secundaria (requer DSSP instalado)
dssp = DSSP(model, "1ubq.pdb")
ss_counts = {'H': 0, 'E': 0, 'C': 0}
for key in dssp.keys():
    ss = dssp[key][2]
    if ss in ss_counts:
        ss_counts[ss] += 1
    else:
        ss_counts['C'] += 1
total = sum(ss_counts.values())
print(f"Alfa-helice: {100*ss_counts['H']/total:.1f}%")
print(f"Beta-folha:  {100*ss_counts['E']/total:.1f}%")
print(f"Loop/coil:   {100*ss_counts['C']/total:.1f}%")
\end{lstlisting}


\section{Cinética Enzimática Avançada}

\subsection{Michaelis-Menten e suas Linearizações}

A equação de Michaelis-Menten descreve a velocidade de reação catalisada por uma enzima
em função da concentração de substrato, derivada do mecanismo mínimo
$E + S \rightleftharpoons ES \to E + P$:

\begin{equation}
v = \frac{V_{\max}[S]}{K_M + [S]}, \quad K_M = \frac{k_{-1} + k_{cat}}{k_1}, \quad V_{\max} = k_{cat}[E]_0
\end{equation}

O $K_M$ é a concentração de substrato que produz metade da velocidade máxima: uma medida
da afinidade aparente da enzima pelo substrato (menor $K_M$ = maior afinidade). O número
de \emph{turnover} $k_{cat}$ é o número de moléculas de produto formadas por segundo por
molécula de enzima saturada. A \emph{eficiência catalítica} $k_{cat}/K_M$ mede a velocidade
de reação a baixas concentrações de substrato e é limitada superiormente pela taxa de
difusão ($10^8$--$10^9$ M$^{-1}$s$^{-1}$); enzimas cataliticamente perfeitas como a
anidrase carbônica e a catalase operam próximas a esse limite.

Historicamente, os parâmetros foram determinados por linearizações. O \emph{gráfico de
Lineweaver-Burk} (duplo recíproco) inverte a equação:
$1/v = (K_M/V_{\max})(1/[S]) + 1/V_{\max}$, produzindo uma reta com intercepto-$y =
1/V_{\max}$ e intercepto-$x = -1/K_M$. Apesar de sua elegância pedagógica, o método
amplifica os erros a concentrações baixas de substrato (grandes valores de $1/[S]$) e
distorce a distribuição de erros experimental. Hoje, o ajuste não-linear direto por
mínimos quadrados (via \texttt{scipy.optimize.curve\_fit}, que usa o algoritmo de
Levenberg-Marquardt internamente) é o método preferido.

O \emph{gráfico de Eadie-Hofstee} ($v$ vs.\ $v/[S]$, com inclinação $-K_M$ e
intercepto-$y$ em $V_{\max}$) distribui os erros de forma mais uniforme e é
particularmente útil para detectar desvios da cinética michaeliana, como cooperatividade.

\subsection{Cooperatividade e o Modelo de Hill}

Em enzimas oligoméricas, a ligação de substrato a um sítio pode alterar a afinidade dos
demais sítios --- \emph{cooperatividade}. A equação de Hill generaliza Michaelis-Menten:

\begin{equation}
v = \frac{V_{\max}[S]^n}{K_d^n + [S]^n}
\end{equation}

onde $n$ é o coeficiente de Hill: $n > 1$ indica cooperatividade positiva (cada ligação
facilita as seguintes); $n < 1$ indica cooperatividade negativa. A hemoglobina, com $n
\approx 2.8$ para ligação de O$_2$, é o exemplo clássico: a curva sigmoidal de saturação
em função da pressão parcial de oxigênio permite que ela carregue O$_2$ eficientemente
nos pulmões (alta $p_{O_2}$) e o libere nos tecidos (baixa $p_{O_2}$).

A cooperatividade positiva produz uma resposta \emph{switch-like}: pequenas variações na
concentração de substrato ao redor de $K_d$ causam grandes variações em $v$. Isso é
funcionalmente relevante para enzimas que regulam fluxos metabólicos ou que atuam como
interruptores moleculares.

Dois modelos mecanísticos explicam a cooperatividade. O \emph{modelo MWC} (Monod-Wyman-Changeux,
1965) postula que a enzima existe em dois estados --- T (tenso, baixa afinidade) e R
(relaxado, alta afinidade) --- e que a transição entre eles é \emph{concertada}: todas as
subunidades mudam de estado simultaneamente. O \emph{modelo KNF} (Koshland-Némethy-Filmer,
1966) propõe transições \emph{sequenciais}: a ligação do substrato a uma subunidade induz
mudança conformacional que afeta apenas as subunidades adjacentes, permitindo estados
híbridos T/R.

\subsection{Tipos de Inibição Enzimática}

Um inibidor enzimático é qualquer molécula que reduz a velocidade catalítica. Quatro
mecanismos principais são distinguidos pela forma como afetam $K_M$ e $V_{\max}$:

Na \emph{inibição competitiva}, o inibidor $I$ compete com o substrato pelo sítio ativo,
ligando-se a $E$ mas não a $ES$. O equilíbrio $E + I \rightleftharpoons EI$ (constante
$K_I$) reduz a fração de enzima disponível para o substrato. Resultado: $V_{\max}$ não
muda (excesso de substrato supera o inibidor) mas $K_M^{app}$ aumenta:
\begin{equation}
v = \frac{V_{\max}[S]}{K_M\!\left(1 + \frac{[I]}{K_I}\right) + [S]}
\end{equation}

Na \emph{inibição não-competitiva}, o inibidor liga ao sítio alostérico em $E$ ou $ES$
com a mesma afinidade ($K_I = K_I'$). Não compete com o substrato, portanto $K_M$ não
muda, mas $V_{\max}^{app}$ diminui proporcionalmente a $1/(1+[I]/K_I)$.

Na \emph{inibição incompetitiva}, o inibidor liga apenas ao complexo $ES$, nunca à enzima
livre. Tanto $K_M$ quanto $V_{\max}$ diminuem pelo mesmo fator: as retas no gráfico de
Lineweaver-Burk são paralelas.

A \emph{inibição mista} (forma geral) ocorre quando o inibidor liga tanto a $E$ quanto a
$ES$, mas com afinidades diferentes ($K_I \neq K_I'$).

O $IC_{50}$ --- concentração de inibidor que reduz a atividade em 50\% --- é a grandeza
mais reportada na descoberta de fármacos. Para inibidores competitivos, a relação de Cheng-Prusoff
conecta $IC_{50}$ e $K_I$: $K_I = IC_{50}/(1 + [S]/K_M)$.

\begin{lstlisting}[language=Python, caption=Ajuste de curva dose-resposta e IC50]
import numpy as np
from scipy.optimize import curve_fit
import matplotlib.pyplot as plt

def michaelis_menten(S, Vmax, Km):
    return Vmax * S / (Km + S)

def dose_response(I, IC50, n):
    return 100 / (1 + (I / IC50)**n)

# Dados cineticos
S = np.array([0.5, 1, 2, 5, 10, 20, 50])   # mM
v = np.array([15, 25, 38, 60, 75, 85, 95])  # umol/min

p_mm, cov_mm = curve_fit(michaelis_menten, S, v, p0=[100, 5])
Vmax, Km = p_mm
print(f"Vmax = {Vmax:.1f} umol/min | Km = {Km:.2f} mM")
print(f"kcat/Km = {Vmax/Km:.2f}  (unidades arbitrarias)")

# Curva dose-resposta do inibidor
I = np.array([0.01, 0.1, 1, 10, 100, 1000])   # uM
act = np.array([98, 90, 70, 30, 8, 2])          # % atividade
p_dr, _ = curve_fit(dose_response, I, act, p0=[10, 1])
IC50, n_hill = p_dr

# Cheng-Prusoff (inibidor competitivo)
S_assay, Km_enz = 5, 2   # mM
Ki = IC50 / (1 + S_assay / Km_enz)
print(f"IC50 = {IC50:.2f} uM | n = {n_hill:.2f}")
print(f"Ki (Cheng-Prusoff) = {Ki:.2f} uM")

# Plot
fig, axes = plt.subplots(1, 2, figsize=(12, 5))
S_fit = np.linspace(0, 50, 200)
axes[0].plot(S, v, 'o', label='Dados'); axes[0].plot(S_fit, michaelis_menten(S_fit, *p_mm))
axes[0].set_xlabel('[S] (mM)'); axes[0].set_ylabel('v (umol/min)'); axes[0].set_title('Michaelis-Menten')
I_fit = np.logspace(-2, 3, 200)
axes[1].semilogx(I, act, 'o'); axes[1].semilogx(I_fit, dose_response(I_fit, *p_dr))
axes[1].axhline(50, color='r', linestyle='--'); axes[1].set_title('Curva Dose-Resposta')
plt.tight_layout(); plt.show()
\end{lstlisting}

\subsection{Ultrassensibilidade de Goldbeter-Koshland}

Em 1981, Goldbeter e Koshland demonstraram que um ciclo de modificação covalente
reversível --- como fosforilação/desfosforilação --- pode produzir resposta
ultrassensível (switch-like) mesmo sem cooperatividade intrínseca nas enzimas, quando
ambas operam perto da saturação (\emph{zero-order ultrasensitivity}). Quando a
concentração total de substrato excede muito os $K_M$ de quinase e fosfatase, uma
pequena variação na razão de atividades ($v_1/v_2$) produz uma grande variação na
fração de proteína fosforilada. Esse mecanismo, que opera com coeficiente de Hill efetivo
$n \gg 1$, é biologicamente fundamental: o switch G2/M do ciclo celular, a ativação de
ERK na cascata MAPK e o switch de secreção de insulina nas células $\beta$ do pâncreas
são todos exemplos de sistemas que exploram a ultrassensibilidade de zero-order.


\section{Cascatas de Sinalização Celular}

\subsection{Princípios de Transdução de Sinal}

\begin{definicaoenv}[Transdução de Sinal]
Transdução de sinal é o processo pelo qual uma célula converte um sinal extracelular ---
hormônio, fator de crescimento, neurotransmissor, força mecânica --- em uma resposta
intracelular específica, por meio de uma cascata sequencial de interações moleculares.
\end{definicaoenv}

O fluxo geral é: ligante extracelular → receptor de membrana → transdutores
intracelulares → efetores → resposta biológica (alteração de expressão gênica,
metabolismo, morfologia ou divisão celular). Cada etapa oferece oportunidade de
amplificação, integração de múltiplos sinais e regulação por feedback.

Os \emph{receptores de membrana} são classificados em quatro famílias principais. Os
\emph{receptores acoplados a proteínas G} (GPCRs) possuem 7 hélices transmembrana e
representam a maior família de receptores humanos ($\sim$800 genes), alvos de 30--40\%
dos fármacos aprovados. Quando ativados por ligante, os GPCRs catalisam a troca de GDP
por GTP na subunidade G$\alpha$, que se dissocia do dímero G$\beta\gamma$ e ativa ou
inibe enzimas efetoras: G$\alpha_s$ ativa a adenilil ciclase (produz cAMP); G$\alpha_i$
a inibe; G$\alpha_q$ ativa a fosfolipase C (produz IP$_3$ e DAG).

Os \emph{receptores tirosina quinase} (RTKs) têm uma única hélice transmembrana e
atividade quinase intrínseca no domínio intracelular. A ligação do fator de crescimento
induz dimerização do receptor, levando à autofosforilação em resíduos de tirosina que
criam sítios de docking para proteínas adaptadoras (Grb2, SOS) e efetoras (PI3K, PLCγ).
Exemplos: EGFR (receptor do fator de crescimento epidermal), VEGFR (fator de crescimento
vascular) e o receptor de insulina.

Os \emph{canais iônicos controlados por ligante} (receptores ionotrópicos) abrem ou
fecham em resposta à ligação do ligante, permitindo fluxo iônico rápido (escala de ms).
O receptor nicotínico de acetilcolina, que medeia a transmissão neuromuscular, é o
paradigma dessa família. Os \emph{receptores nucleares} são fatores de transcrição
intracelulares que, ao ligar hormônios lipofílicos (esteroides, hormônios tireoidianos,
retinoides), se ativam e regulam diretamente a transcrição de genes alvo.

\subsection{A Cascata MAPK e Amplificação}

A via MAPK (Mitogen-Activated Protein Kinase) é um módulo de três quinases em série ---
MAPKKK → MAPKK → MAPK --- conservado desde leveduras a humanos. Na via ERK clássica em
mamíferos: fator de crescimento ativa RTK → Grb2/SOS recrutado → Ras ativado (GDP → GTP)
→ Raf (MAPKKK) ativado → MEK (MAPKK) duplamente fosforilado → ERK (MAPK) fosforilado.
ERK ativa transfatorea de transcrição (c-Fos, Elk-1) e quinases citoplásmicas,
coordenando proliferação, diferenciação e sobrevivência celular.

A cascata opera como amplificador exponencial: uma única RTK ativa múltiplos Ras, cada
Ras ativa múltiplos Raf, etc. Se cada quinase ativa 10 moléculas da próxima, a cascata
de três etapas produz amplificação de $10^3$. Ao mesmo tempo, as quinases e suas
fosfatases oponentes criam circuitos de feedback --- incluindo feedback negativo (ERK
fosforila e inativa SOS, limitando a duração do sinal) e feedback positivo (ERK
fosforila Raf em certas condições, criando bistabilidade) --- que controlam a forma
temporal e a duração da resposta.

\begin{lstlisting}[language=Python, caption=Modelo simplificado de cascata MAPK]
import numpy as np
from scipy.integrate import odeint
import matplotlib.pyplot as plt

def mapk_cascade(y, t, k_on, k_off, phos):
    MAPKKK, MAPKK, MAPK = y
    # Ativacao por sinal externo (step function em t=5)
    signal = 1.0 if t >= 5 else 0.0
    dMAPKKK = k_on * signal - k_off * MAPKKK
    dMAPKK  = phos * MAPKKK - k_off * MAPKK
    dMAPK   = phos * MAPKK  - k_off * MAPK
    return [dMAPKKK, dMAPKK, dMAPK]

t = np.linspace(0, 50, 1000)
sol = odeint(mapk_cascade, [0, 0, 0], t, args=(0.5, 0.1, 0.8))

plt.figure(figsize=(9, 5))
for i, label in enumerate(['MAPKKK*', 'MAPKK*', 'MAPK*']):
    plt.plot(t, sol[:, i], label=label, linewidth=2)
plt.axvline(5, color='gray', linestyle='--', label='Sinal ON')
plt.xlabel('Tempo (min)'); plt.ylabel('Concentracao (uM)')
plt.title('Cascata MAPK: Amplificacao e Atraso')
plt.legend(); plt.grid(True); plt.tight_layout(); plt.show()
\end{lstlisting}

\subsection{Segundos Mensageiros}

Os \emph{segundos mensageiros} são moléculas difusíveis que propagam e amplificam o sinal
do receptor para efetores distantes no citoplasma. O \emph{AMP cíclico} (cAMP) é
produzido pela adenilil ciclase ativada por G$\alpha_s$ e ativa a proteína quinase A (PKA),
que fosforila dezenas de substratos incluindo o fator de transcrição CREB. As
\emph{fosfodiesterases} hidrolisam cAMP, encerrando o sinal; inibidores de
fosfodiesterases (como a cafeína e o sildenafil) amplificam a sinalização por cAMP e cGMP.

O \emph{cálcio intracelular} ($[\text{Ca}^{2+}]_i$) é mantido em $\sim$100 nM no repouso
--- quatro ordens de magnitude abaixo da concentração extracelular ($\sim$1 mM). Sinais
que ativam fosfolipase C levam à produção de IP$_3$, que abre canais de liberação de
cálcio no retículo endoplasmático, elevando $[\text{Ca}^{2+}]_i$ para $\sim$1 µM. O
cálcio livre liga-se à calmodulina, ativando quinases CaM-dependentes (CaMKII) e PKC, com
consequências para contração muscular, exocitose, apoptose e plasticidade sináptica. As
oscilações de cálcio discutidas no Capítulo 4 são um exemplo de como a dinâmica de um
segundo mensageiro carrega informação na forma temporal, não apenas na amplitude.


\section{Modificações Pós-Traducionais}

\begin{definicaoenv}[Modificação Pós-Traducional (PTM)]
Modificação química covalente de uma proteína que ocorre após a tradução ribosomal.
PTMs expandem enormemente a diversidade funcional do proteoma: os $\sim$20.000 genes
humanos geram mais de 1.000.000 de \emph{proteoformas} distintas quando se contam
todas as isoformas de splicing e combinações de PTMs.
\end{definicaoenv}

Mais de 400 tipos de PTMs são conhecidas, mas quatro dominam a literatura de sinalização
e regulação: fosforilação, acetilação, metilação e ubiquitinação.

\subsection{Fosforilação}

A fosforilação é a PTM mais dinâmica e estudada. A adição de um grupo fosfato (PO$_4^{3-}$)
a resíduos de serina (90\%), treonina (10\%) ou tirosina ($<$1\%) é catalisada por
\emph{quinases}; a remoção é catalisada por \emph{fosfatases}. O genoma humano codifica
$\sim$518 quinases (o \emph{kinome}) --- uma das maiores famílias proteicas --- e $\sim$150
fosfatases. O estado de fosforilação de uma proteína num dado momento reflete o equilíbrio
dinâmico entre a atividade dessas duas classes.

A fosforilação altera a carga local da proteína (de 0 para $-2$ a pH fisiológico) e pode
induzir mudanças conformacionais, criar ou destruir sítios de ligação para domínios
proteicos (SH2 reconhece fosfotirosina; 14-3-3 reconhece fosfoserina/treonina), ou
desencadear localização subcelular alterada e degradação proteossomal. O estado de
fosforilação de centenas de proteínas simultaneamente constitui o \emph{fosfoproteoma}
de uma célula, que muda em segundos a minutos em resposta a sinais.

\subsection{Acetilação e Metilação: o Código de Histonas}

A acetilação de resíduos de lisina (adição de CH$_3$CO) neutraliza a carga positiva do
grupo $\varepsilon$-amino, reduzindo a interação eletrostática com o DNA carregado
negativamente. Nas histonas, a hiperacetilação (especialmente H3K27ac e H4K16ac) relaxa
a cromatina e ativa a transcrição; a hipoacetilação, promovida pelas histona deacetilases
(HDACs), compacta a cromatina e silencia genes. Inibidores de HDAC são fármacos
anticancerígenos aprovados (vorinostat, romidepsin).

A metilação de lisinas é contextualmente mais sutil: a mesma modificação em diferentes
sítios pode ter efeitos opostos. H3K4me3 (trimetilação de K4 da histona H3) marca
promotores de genes ativamente transcritos; H3K9me3 e H3K27me3 marcam cromatina
reprimida e heterocromatina constitutiva, respectivamente. A combinação de modificações
em múltiplos resíduos de histona --- o \emph{código de histonas} (Strahl e Allis, 2000)
--- é reconhecida por proteínas leitoras (\emph{readers}) que recrutam complexos de
remodelamento de cromatina.

\subsection{Ubiquitinação e a Via do Proteassoma}

A ubiquitina é uma proteína de 76 aminoácidos altamente conservada que é covalentemente
ligada a resíduos de lisina de proteínas-alvo num processo de três etapas: ativação por
E1 (enzima ativadora, 2 humanas), transferência a E2 (enzima conjugadora, $\sim$40
humanas) e ligação ao substrato mediada por E3 (ubiquitina ligase, $>$600 humanas). A
especificidade do sistema é conferida pelas E3 ligases, que reconhecem sequências
degron específicas --- frequentemente geradas por fosforilação prévia (fosfo-degron).

A cadeia de ubiquitina montada no substrato tem significados distintos segundo o resíduo
de lisina interno que conecta as ubiquitinas: K48-Ub (cadeia ligada em K48) é o sinal
clássico de degradação pelo proteassoma 26S; K63-Ub sinaliza para vias não-degradativas
(reparo de DNA, endocitose, sinalização inflamatória). A mono-ubiquitinação (uma única
ubiquitina) regula endocitose de receptores e translesão de DNA.

\subsection{Crosstalk entre PTMs}

PTMs não atuam isoladamente --- elas se regulam mutuamente em redes complexas de
\emph{crosstalk}. A fosforilação de um resíduo de serina adjacente a uma lisina pode
bloquear a acetilação dessa lisina por impedimento estérico ou carga. A ubiquitinação de
Lys$_{382}$ de p53 compete com a acetilação do mesmo resíduo: a acetilação ativa p53
(estabilizando-o e aumentando sua afinidade pelo DNA), enquanto a ubiquitinação leva à
sua degradação. A O-GlcNAcilação compete com a fosforilação pelos mesmos resíduos de
serina e treonina, com implicações no controle de glicose e diabetes. O crosstalk entre
PTMs nas histonas coordena o estado transcricional de regiões cromossômicas inteiras.

A detecção de PTMs em escala proteômica é feita por espectrometria de massa (LC-MS/MS):
peptídeos trípticos de proteínas digeridas são separados por cromatografia líquida, e os
espectros de fragmentação identificam a sequência e a massa exata de cada modificação.
As bases de dados PhosphoSitePlus (phosphosite.org) e dbPTM integram dados experimentais
de PTMs de múltiplas espécies.


\section{Proteômica}

\begin{definicaoenv}[Proteoma]
O proteoma é o conjunto completo de proteínas expressas por uma célula, tecido ou
organismo num determinado momento e condição. Ao contrário do genoma, que é relativamente
estável, o proteoma é altamente dinâmico: muda com o tipo celular, o estágio de
desenvolvimento, os sinais extracelulares, o estado metabólico e o estresse.
\end{definicaoenv}

A distinção entre transcriptoma e proteoma é biologicamente fundamental: a correlação entre
mRNA e proteína para o mesmo gene é tipicamente $r^2 \approx 0.4$--$0.6$, dependendo do
tecido e das condições. As razões incluem diferentes taxas de tradução por mRNA, variações
nas taxas de degradação proteica, regulação pós-transcricional por miRNAs, e a multiplicação
de proteoformas por splicing alternativo e PTMs. O proteoma humano, apesar de derivar de
$\sim$20.000 genes, provavelmente compreende mais de 1.000.000 de proteoformas distintas.

\subsection{Espectrometria de Massa em Proteômica}

A espectrometria de massa (MS) é a tecnologia central da proteômica. No fluxo de trabalho
\emph{bottom-up} (shotgun proteomics), proteínas são primeiro digeridas com a protease
tripsina (que cliva após arginina e lisina), produzindo peptídeos de 6--25 aminoácidos.
Esses peptídeos são separados por cromatografia líquida de ultra-alta resolução (UHPLC)
e ionizados por electrospray (ESI), gerando íons com múltiplas cargas que são analisados
pelo espectrômetro.

Em modo DDA (\textit{data-dependent acquisition}), o espectrômetro seleciona automaticamente
os íons mais abundantes do espectro de massa de survey (MS1) para fragmentação por colisão
(CID ou HCD), gerando espectros de fragmentação (MS2) que revelam a sequência do peptídeo
pela série de íons $b$ e $y$. A busca desses espectros contra bancos de dados de sequências
(UniProt, RefSeq) por softwares como Mascot, Sequest ou MaxQuant identifica e quantifica as
proteínas presentes na amostra.

\subsection{Quantificação Proteômica}

Três estratégias principais de quantificação são empregadas. A abordagem \emph{label-free}
usa a intensidade dos picos no espectro MS1 (LFQ, \textit{label-free quantification}) ou
a contagem de espectros de MS2 para estimar a abundância relativa. É simples e aplicável a
qualquer tipo de amostra, mas tem variabilidade inter-amostra maior.

O \emph{SILAC} (\textit{Stable Isotope Labeling by Amino Acids in Cell Culture}) incorpora
metabolicamente aminoácidos marcados com isótopos pesados ($^{13}$C, $^{15}$N) nas células
experimentais, enquanto o controle usa aminoácidos leves. As amostras são misturadas antes
da digestão, e a razão dos pares de picos leve/pesado fornece quantificação altamente
precisa --- tipicamente coeficiente de variação $<$10\%.

A marcação química por reagentes \emph{TMT} (\textit{Tandem Mass Tags}) ou \emph{iTRAQ}
permite multiplexar 6--18 amostras num único experimento. Os reagentes TMT têm a mesma
massa (isobáricos), mas liberam íons repórteres de massas distintas na fragmentação MS2,
permitindo quantificação simultânea de múltiplas condições.

\begin{lstlisting}[language=Python, caption=Processamento de dados de proteomica LFQ]
import pandas as pd
import numpy as np
from scipy.stats import ttest_ind
from statsmodels.stats.multitest import multipletests

# Simular dados LFQ (log2-transformados)
np.random.seed(42)
n_prot = 1000
ctrl  = np.random.normal(20, 2, (n_prot, 3))
treat = np.random.normal(20, 2, (n_prot, 3))
# 50 proteinas super-expressas no tratamento
treat[:50] += np.random.normal(2, 0.5, (50, 3))

# Calcular fold change e p-valor
log2fc = treat.mean(1) - ctrl.mean(1)
p_vals = np.array([ttest_ind(ctrl[i], treat[i]).pvalue for i in range(n_prot)])

# Correcao Benjamini-Hochberg
_, p_adj, _, _ = multipletests(p_vals, method='fdr_bh')
sig = (p_adj < 0.05) & (np.abs(log2fc) > 1)
print(f"Proteinas significativas: {sig.sum()}")
print(f"Super-expressas: {((log2fc > 1) & (p_adj < 0.05)).sum()}")
print(f"Sub-expressas:   {((log2fc < -1) & (p_adj < 0.05)).sum()}")
\end{lstlisting}

\subsection{Interações Proteína-Proteína}

O interatoma proteico --- o conjunto completo de interações físicas entre proteínas ---
é uma das redes biológicas mais estudadas. Métodos experimentais complementares
mapeiam diferentes aspectos: o Y2H (\textit{yeast two-hybrid}) detecta interações
binárias entre pares de proteínas em alto-throughput (mas tem alta taxa de falsos
positivos); o AP-MS (\textit{affinity purification-mass spectrometry}) purifica
complexos nativos marcando uma proteína isca com uma tag de afinidade e identificando
os parceiros copurificados por espectrometria de massa (maior fidelidade, mas perde
interações transitórias); o BioID e o APEX utilizam enzimas de biotinilação
(BirA*, APEX2) fundidas à proteína de interesse para marcar quimicamente proteínas
em sua vizinhança imediata in vivo --- ideal para interações transitórias e proteínas
de membrana inacessíveis ao AP-MS.

As bases de dados STRING (string-db.org), BioGRID e IntAct integram esses dados para
construir o interatoma de referência de múltiplos organismos. A análise de redes de
interação proteica, cobrindo tópicos como identificação de hubs, módulos funcionais e
proteínas essenciais, foi discutida em detalhes no Capítulo 3.


\section{Exercícios}

\begin{exercicio}
Uma enzima tem $K_M = 5$ mM e $V_{\max} = 100$ $\mu$mol/min. (a) Calcule a velocidade
quando $[S] = 15$ mM. (b) Determine a concentração de substrato necessária para atingir
$v = 75$ $\mu$mol/min. (c) Se a concentração total de enzima é $[E]_0 = 1$ $\mu$M,
calcule $k_{cat}$ e $k_{cat}/K_M$. Compare com o limite de difusão.
\end{exercicio}

\begin{exercicio}
Um inibidor competitivo tem $K_I = 2$ mM. Se $K_M = 5$ mM e $V_{\max} = 100$ $\mu$mol/min:
(a) calcule o $K_M^{app}$ quando $[I] = 10$ mM; (b) construa o gráfico de Lineweaver-Burk
com e sem inibidor; (c) que concentração de substrato é necessária para atingir 90\% de
$V_{\max}$ na presença do inibidor?
\end{exercicio}

\begin{exercicio}[Hill]
Para uma enzima cooperativa com $n = 3$ e $K_d = 8$ $\mu$M: (a) calcule $v/V_{\max}$
para $[S] = 4$, 8 e 16 $\mu$M; (b) compare com a hipérbole de Michaelis-Menten
($n = 1$) para os mesmos valores; (c) construa o gráfico de Hill linearizado
($\log[v/(V_{\max}-v)]$ vs.\ $\log[S]$) e interprete a inclinação.
\end{exercicio}

\begin{exercicio}
Explique por que: (a) Glicina tem maior liberdade conformacional que outros aminoácidos
no diagrama de Ramachandran; (b) Prolina tem ângulo $\phi$ restrito; (c) proteínas
intrinsecamente desordenadas podem ser funcionalmente ativas apesar de não terem estrutura
tridimensional fixa.
\end{exercicio}

\begin{exercicio}[MAPK]
Na cascata MAPK, cada quinase ativa em média 10 moléculas da quinase seguinte. (a) Quantas
moléculas de ERK são ativadas por 1 RTK? (b) Se há 100 RTKs ativados na membrana, quantas
ERK fosforiladas se espera? (c) Como um feedback negativo de ERK sobre SOS (com constante
de meia-vida de inativação de 5 min) afetaria a duração da resposta? Escreva as EDOs
para o sistema com feedback.
\end{exercicio}

\begin{exercicio}[PTMs]
Uma proteína tem Lys$_{48}$ (pode ser acetilada OU ubiquitinada) e Ser$_{132}$ (pode ser
fosforilada OU O-GlcNAcilada), além de estar em sua forma não-modificada. (a) Quantas
proteoformas distintas são possíveis? (b) Qual combinação de PTMs deveria ativar a proteína
(acetilação ativa, fosforilação ativa) versus promover sua degradação (ubiquitinação K48)?
(c) Que técnica experimental usaria para detectar cada PTM individualmente?
\end{exercicio}

\begin{exercicio}[BioPython]
Usando BioPython: (a) baixe a estrutura PDB da lisozima de galinha (PDB: 1LZT); (b) calcule
a proporção de resíduos em $\alpha$-hélice, $\beta$-folha e loop; (c) identifique os quatro
resíduos de cisteína que formam pontes dissulfeto e calcule a distância entre os átomos S$\gamma$
de cada par; (d) liste todos os resíduos com menos de 5 Å do ligante NAG (N-acetilglucosamina).
\end{exercicio}

\begin{exercicio}[Proteômica --- Projeto]
Acesse o banco ProteomeXchange e baixe dados de proteômica quantitativa (label-free ou TMT)
de um experimento de sua escolha. Construa um pipeline completo: (a) log2-transformação e
normalização (mediana ou quantílica); (b) filtro de valores ausentes ($>$70\% de valores
válidos por grupo); (c) teste $t$ por proteína com correção FDR de Benjamini-Hochberg;
(d) volcano plot; (e) análise de enriquecimento GO/KEGG das proteínas diferencialmente
expressas; (f) rede de interações das proteínas significativas via STRING. Interprete
biologicamente os resultados.
\end{exercicio}

\section{Notas Bibliográficas}

Para estrutura de proteínas e bioquímica geral, Berg, Tymoczko e Stryer \cite{berg2019}
(\textit{Biochemistry}, 9ª ed.) é a referência textual mais utilizada. O diagrama de
Ramachandran foi introduzido em Ramachandran et al.\ (1963) \cite{ramachandran1963}. As
proteínas intrinsecamente desordenadas são revisadas em Dyson e Wright (2005)
\cite{dyson2005}.

Para cinética enzimática, Cornish-Bowden \cite{cornish2012} (\textit{Fundamentals of
Enzyme Kinetics}) é a referência matemática mais completa. O modelo MWC foi publicado em
Monod, Wyman e Changeux (1965) \cite{mwc1965}; o modelo KNF em Koshland, Némethy e Filmer
(1966) \cite{knf1966}. A ultrassensibilidade de zero-order foi demonstrada por Goldbeter e
Koshland (1981) \cite{goldbeter1981}.

Para sinalização celular, Lodish et al.\ \cite{lodish2021} (\textit{Molecular Cell
Biology}) cobre os mecanismos com rigor molecular. A descoberta dos GPCRs e das proteínas
G foi premiada com o Nobel de 2012 (Lefkowitz e Kobilka). Para PTMs, Mann e Jensen (2003)
\cite{mann2003} introduziram a análise de PTMs por MS, e Venne et al.\ (2014)
\cite{venne2014} revisam o crosstalk. O conceito de código de histonas foi proposto por
Strahl e Allis (2000) \cite{strahl2000}.

Para proteômica, Aebersold e Mann (2016) \cite{aebersold2016} é a revisão mais abrangente
sobre exploração proteômica por MS. O software MaxQuant (Cox e Mann, 2008 \cite{cox2008})
é o padrão para análise de dados de MS proteômico.

\begin{thebibliography}{99}

\bibitem{berg2019}
Berg J.M., Tymoczko J.L., Stryer L. (2019).
\textit{Biochemistry}, 9\textordfeminine\ ed.
W.H. Freeman.

\bibitem{ramachandran1963}
Ramachandran G.N., Ramakrishnan C., Sasisekharan V. (1963).
Stereochemistry of polypeptide chain configurations.
\textit{Journal of Molecular Biology} \textbf{7}: 95--99.

\bibitem{dyson2005}
Dyson H.J., Wright P.E. (2005).
Intrinsically unstructured proteins and their functions.
\textit{Nature Reviews Molecular Cell Biology} \textbf{6}: 197--208.

\bibitem{cornish2012}
Cornish-Bowden A. (2012).
\textit{Fundamentals of Enzyme Kinetics}, 4\textordfeminine\ ed.
Wiley-Blackwell.

\bibitem{mwc1965}
Monod J., Wyman J., Changeux J.P. (1965).
On the nature of allosteric transitions: a plausible model.
\textit{Journal of Molecular Biology} \textbf{12}: 88--118.

\bibitem{knf1966}
Koshland D.E., Nemethy G., Filmer D. (1966).
Comparison of experimental binding data and theoretical models in proteins containing
subunits.
\textit{Biochemistry} \textbf{5}: 365--385.

\bibitem{goldbeter1981}
Goldbeter A., Koshland D.E. (1981).
An amplified sensitivity arising from covalent modification in biological systems.
\textit{Proceedings of the National Academy of Sciences} \textbf{78}: 6840--6844.

\bibitem{lodish2021}
Lodish H. et al. (2021).
\textit{Molecular Cell Biology}, 9\textordfeminine\ ed.
W.H. Freeman.

\bibitem{mann2003}
Mann M., Jensen O.N. (2003).
Proteomic analysis of post-translational modifications.
\textit{Nature Biotechnology} \textbf{21}: 255--261.

\bibitem{venne2014}
Venne A.S. et al. (2014).
The next level of complexity: Crosstalk of posttranslational modifications.
\textit{Proteomics} \textbf{14}: 513--524.

\bibitem{strahl2000}
Strahl B.D., Allis C.D. (2000).
The language of covalent histone modifications.
\textit{Nature} \textbf{403}: 41--45.

\bibitem{aebersold2016}
Aebersold R., Mann M. (2016).
Mass-spectrometric exploration of proteome structure and function.
\textit{Nature} \textbf{537}: 347--355.

\bibitem{cox2008}
Cox J., Mann M. (2008).
MaxQuant enables high peptide identification rates, individualized p.p.b.-range mass
accuracies and proteome-wide protein quantification.
\textit{Nature Biotechnology} \textbf{26}: 1367--1372.

\end{thebibliography}
